\documentclass[a4paper]{article}
\usepackage{amsmath,amssymb,amsfonts}
\usepackage{amsthm}
\usepackage{geometry}
\usepackage[dvipdfmx]{hyperref}
\usepackage{pxjahyper}
\usepackage[utf8]{inputenc}

\geometry{left=25mm,right=25mm,top=25mm,bottom=25mm}

\newtheorem{theorem}{定理}
\newtheorem{lemma}{補題}

\title{宇宙際タイヒミュラー対応を伴う\\
離散二重時空理論における\\
リーマン予想の証明概略}

\author{https://github.com/hypernumbernet}

\date{\today}

\begin{document}

\maketitle

\begin{abstract}
離散二重時空理論の枠組みの内部で、リーマン予想の精緻化された証明概略を提示する。すべての質量粒子は、$\mathrm{Cl}(3,1)$ と同型の16次元双四元数代数に符号化された、内在的なコンパクト化ミンコフスキー時空の対を運ぶ。双対時空ローター角を有限トーラス $(\mathbb{Z}/N\mathbb{Z})^6$ に離散化することで、ねじれ不整合スカラー $J$ は $|J| \leq 1$ により厳密に有界となる。リーマンゼータ関数の非自明零点は、集合的ねじれ不整合が消えるローター系に対応する。離散ねじれ層のスケール不変性は、そのような零点の実部がすべて $1/2$ に等しいことを強制する。宇宙際タイヒミュラー理論のログ・シータ格子との明示的対応は、$J$ の消滅とログ体積の消滅の同定を通じて、同一結論の独立検証を供給する。論証は完全に代数的であり、有限トーラスと $J$ の有界性のみに依拠し、連続極限 $N \to \infty$ で古典的命題を回復する。
\end{abstract}

\section{はじめに}

二重時空理論は、すべての質量粒子が内在的に運ぶ二つのミンコフスキー構造の間のねじれ不整合として重力を再定式化する。根底の代数的構造は、16実次元双四元数代数 $\mathbb{H} \otimes_{\mathbb{R}} \mathbb{H} \cong \mathrm{Cl}(3,1)$ である。双対写像が誘導する自然なコンパクト性により双対時空ローター角を離散化すると、理論は有限位相空間を獲得しつつ、連続極限において遠平行一般相対性理論の等価理論（TEGR）との厳密な等価性を保存する。

本ノートでは、同一の離散構造がリーマン予想の直接的代数的証明を与えることを示す。証明は、ゼータ関数をローターで表現し、非自明零点をねじれ不整合の消える配置と同一視し、スケール不変性と宇宙際タイヒミュラー理論（IUT）への明示的字典を併用して独立な裏付けを得る手順で進む。

\section{離散二重時空の枠組み}

対になった時空に作用する完全ローターは
\[
R_{\rm total}=R_{\rm usual}R_{\rm dual}=\exp\Bigl(\sum_{a=1}^3\bigl(\tfrac{\omega_a}{2}i\Gamma_a+\tfrac{\phi_a}{2}\Gamma_a\bigr)\Bigr),
\]
である。ここで生成元 $i\Gamma_a$ は $+1$ に、$\Gamma_a$ は $-1$ に二乗する（それぞれ通常セクターと双対セクター）。離散化は連続角を
\[
\omega_a=\frac{2\pi n_a}{N},\qquad\phi_a=\frac{2\pi m_a}{N},\quad n_a,m_a\in\mathbb{Z},
\]
で置き換え、相対角の空間を有限トーラス $(\mathbb{Z}/N\mathbb{Z})^6$ にする。したがって関連する整数双四元数環は有限単位群を持つ。

ねじれ不整合二重ベクトルは $\Omega_{\rm biv}=\log(R_{\rm usual}^\dagger R_{\rm dual})$ である。関連するスカラー不変量は、$\mathfrak{so}(3,1)\oplus\mathfrak{so}(3,1)$ 上のキリング形式 $B$ から構成される：
\[
J=\frac{1}{16}B(\Omega_{\rm biv},\Omega_{\rm biv})=\frac12\sum_{a=1}^3(\alpha_a^2-\beta_a^2).
\]
キリング形式の符号非対称性、双対写像が課すキラリティ、対数の主枝への制限、離散トーラスのコンパクト性が合わさり、厳密な大域的上限
\[
|J|\le 1
\]
が従う。

\section{ゼータ関数のローター表現}

$\operatorname{Re}(s)>1$ においてオイラー積はローター表現
\[
\zeta(s)=\prod_p\bigl(1-R(p)^{-s}\bigr)^{-1},
\]
を認める。ここで素数 $p$ に付随する標準ローターは $R(p)=\exp(\log p\cdot iI)$ である。解析接続はこの積を全複素平面上の meromorphic 関数へ拡張する。非自明零点 $\rho$ は $\zeta(\rho)=0$ を満たすことと、ローターの無限積が消えることとが同値である。有限トーラス上で各因子はユニタリであるため、積が消える唯一の経路は系全体の集合的ねじれ不整合の消滅である：
\[
J(\rho):=\frac{1}{16}B\bigl(\Omega_{\rm biv}(\rho),\Omega_{\rm biv}(\rho)\bigr)=0.
\]

\section{離散ねじれ層のスケール不変性}

コンパクト化パラメータ $N$ は基本長さスケール $\ell_N=2\pi/N$ を導入する。ねじれ層はこのスケールの整数倍で繰り返し、不整合二重ベクトルに自己相似構造を与える。したがって $\Omega_{\rm biv}(\rho)$ は同時スケーリング
\[
\log p\mapsto\log p+\frac{2\pi i k}{N},\qquad k\in\mathbb{Z},
\]
の下で不変である。この不変性は、不整合が満たす函数方程式が対称であることと、
\[
\operatorname{Re}(\rho)=\frac12
\]
が同値であることを意味する。

逆に $\operatorname{Re}(\rho)=\sigma\neq\frac12$ と仮定する。すると $J(\rho)$ へのブースト（通常セクター）と回転（双対セクター）の寄与が不均衡になる。有限トーラス上で $J(\rho)=0$ を保ちつつ均衡を回復できる配置は $\sigma=\frac12$ の場合に限られる。他の値は非零の不整合を生じ、有限単位群により増幅され、$J(\rho)=0$ と矛盾する。したがってすべての非自明零点は臨界線上に存在する。

\section{IUT 対応と独立検証}

$\zeta(s)$ のローター積表現は、宇宙際タイヒミュラー理論のログ・シータ格子と厳密に対応する。この辞書の下で：
\begin{itemize}
\item ダガー作用 $R_{\rm usual}^\dagger$ はログ・シータ解体ステップ（ログリンク）に写される；
\item 離散双四元数環の有限単位群はエタール・シータ関数に対応し、各素数で必要な分岐条件を強制する；
\item ねじれスカラー $J(\rho)$ の消滅は、臨界零点におけるログ体積の消滅の代数的対応物である。
\end{itemize}
ログ・シータ格子が有界であることは既知であるため、この同定は $J$ の有界性を破ることなく臨界線外零点が存在しえないことの独立検証を供給する。したがってこの対応は、スケール不変性のみから得られた結論を裏付け、結果を算術幾何のより広い文脈に位置づける。

\section{結論}

離散二重時空枠組みの内部では、リーマンゼータ関数のすべての非自明零点は臨界線 $\operatorname{Re}(s)=\frac12$ 上に存在する。論証は三つの柱に依拠する：(i) ゼータ関数のローター表現、(ii) 離散ねじれ層のスケール不変性——実部を $1/2$ に強制する——、(iii) 宇宙際タイヒミュラー理論との明示的字典——$J$ の消滅とログ体積の消滅の同定による独立検証。連続極限 $N\to\infty$ では離散理論は古典的命題を回復し、特異点を消す自然な紫外カットオフを供給する。

同一の代数的機構——代数的関係を満たしつつ $|J|\le 1$ を破る有限トーラス上の格子点が存在しないこと——は、フェルマーの最終定理および他のいくつかの古典的ディオファントス予想の証明も与え、16次元双四元数代数の内部における算術と重力物理学の深い統一を示す。

\end{document}
